\startenvironment octcv
\usemodule[zhfonts]
\definecolor[CharmRed][r=.7,g=.4,b=.4]
\setevariables[cv][linethickness=4pt,linecolor=CharmRed]
\setvariables[cvtitle][style=\bfc]
\setvariables[cvsubject][style=\bfa]
\setvariables
  [cvsubject]
  [bottomlinethickness=6pt,bottomlinecolor=CharmRed,splitlinethickness=2pt]
\setuplayout[backspace=.075\paperwidth,
             width=\dimexpr\paperwidth-2\backspace\relax,
             topspace=\backspace,header=0cm,footer=0cm,
             height=\dimexpr\paperheight -2\topspace\relax]
\setuppagenumbering[alternative=singlesided,location=]
\setupheads[indentnext=yes]
\setuphead[title,subject][before={\blank[halfline]},after={\blank[halfline]}]
\setuphead[title][style=\getvariable{cvtitle}{style}]
\setuphead[subject][style=\getvariable{cvsubject}{style}]
\define[1]\setsubjectwidth{%
  \setbox0\hbox{\getvariable{cvsubject}{style} #1}%
  \setvariables[cvsubject][width=\the\wd0]%
}
\setuphead[subject][command=\subjectCMD]
\define[2]\subjectCMD{%
  \setsubjectwidth{#2}\cvframed{#2}%
}
\define[3]\DefSym{%
  \definesymbol
    [#1]
    [{\offset[bottomoffset=#2]{\externalfigure[icons/#1.png][height=#3]}}]
  \expandafter\def\csname cv#1\endcsname##1{\symbol[#1]\ ##1}%
}
\DefSym{birthday}{-.2em}{1.25em} % 生日
\DefSym{sex}{-.1em}{1.2em} % 性别
\DefSym{address}{-.2em}{1.2em} % 地址
\DefSym{education}{-.2em}{1.2em} % 学历
\DefSym{phone}{-.2em}{1.1em} % 电话
\DefSym{mail}{-.2em}{.9em} % 邮箱
\DefSym{im}{-.2em}{1.25em}
\define[8]\cvstage{%
  \starttabulate[|lw(#1\textwidth)|cw(#3\textwidth)|cw(#5\textwidth)|rw(#7\textwidth)|]
                [unit=0cm,before={\blank[halfline]},after={\blank[halfline]}]
  \NC #2 \NC #4 \NC #6 \NC #8 \NC\NR
  \stoptabulate
}
\definesymbol
  [RightArrow]
  [{\offset[bottomoffset=.1em]{\useMPgraphic{RightArrow}}}]
\startuseMPgraphic{RightArrow}
ahvariant := 1; ahdimple := 1; ahlength := 5pt;
drawarrow (0, 0) -- (5mm, 0) withpen pencircle;
\stopuseMPgraphic
\def\RightArrow{\symbol[RightArrow]}
\setupindenting[first,always,2em]
\setupinterlinespace[line=1.5\bodyfontsize]
% width 是列表项左侧符号区域宽度，将其设置为 2em，
% 是为了与段落首行缩进一致。
\setupitemize[1][framesym,nowhite][width=2em]
\definesymbol[framesym][{\offset[bottomoffset=-.1em]{\useMPgraphic{framesym}}}]
\startuseMPgraphic{framesym}
numeric u; path p; color squarecolor;
u := BodyFontSize;
p := fullsquare scaled .625u;
squarecolor := \MPcolor{\getvariable{cvlist}{color}};
fill p shifted (.2u, -.2u) withcolor darkgray;
fill p withcolor squarecolor;
\stopuseMPgraphic
\setvariables[cvlist][color=CharmRed]
\setvariables[photo][x=16cm,y=24.25cm,shapecolor=CharmRed]
\startuseMPgraphic{cvbackground}
draw (0, 0) withcolor background;
draw (PaperWidth, PaperHeight) withcolor background;
numeric linethickness; color linecolor;
linethickness := \getvariable{cv}{linethickness};
linecolor := \MPcolor{\getvariable{cv}{linecolor}};
pair a, b;
a := (BackSpace, PaperHeight - .5TopSpace);
b := (BackSpace + TextWidth, ypart a);
draw a -- b withpen pencircle scaled linethickness withcolor linecolor;
pair c;
a := (BackSpace, .5TopSpace); b := (BackSpace + TextWidth, ypart a);
c := .5[a, b];
drawoptions(withpen pencircle scaled linethickness withcolor linecolor);
if \number\lastuserpage = 1:
  draw a -- b;
else:
  numeric delta; delta := 1.75BodyFontSize;
  picture pn; pn := textext("{\bf\userpagenumber}");
  draw (a -- (c shifted (-delta, 0)))
          .. (c shifted (0, .5delta)) .. ((c shifted (delta, 0)) -- b);
  draw pn shifted c;
fi;
% 载入照片
picture photo;
photo := externalfigure "icons/photo.png";
% 确定照片尺寸
numeric f;
f := (bbwidth photo) / (bbheight photo);
photo := photo xysized (3.5cm, 3.5cm / f);
% 裁剪为圆形区域
path shape; color shapecolor;
shape := unitcircle scaled 3.5cm;
shapecolor := \MPcolor{\getvariable{photo}{shapecolor}};
clip photo to shape;
% 照片定位，需要试探几次，方能找到合适位置
pair location;
location := (\getvariable{photo}{x}, \getvariable{photo}{y});
photo := photo shifted location;
shape := shape shifted location;
% 只在首页显示照片
if \number\userpagenumber = 1:
    draw photo;
    draw shape withpen pencircle scaled 2pt withcolor shapecolor;
fi;
\stopuseMPgraphic
\defineoverlay[cvbackground][\useMPgraphic{cvbackground}]
\setupbackgrounds[page][background=cvbackground]
\defineframed
  [cvframed]
  [width=\textwidth,frame=off,boffset=2pt,
   align={flushleft,lohi},background=BottomLines]
\defineoverlay[BottomLines][\useMPgraphic{BottomLines}]
\startuseMPgraphic{BottomLines}
numeric thi, w; color linecolor;
thi := \getvariable{cvsubject}{bottomlinethickness};
w := \getvariable{cvsubject}{width};
linecolor := \MPcolor{\getvariable{cvsubject}{bottomlinecolor}};
% 画底线
path bottomline;
bottomline := (.5thi, 0) -- (w, 0);
draw bottomline withpen pencircle scaled thi withcolor transparent(1, .5, linecolor);
% 画版心贯穿线
path splitline;
splitline := (0, 0) -- (OverlayWidth, 0);
thi := \getvariable{cvsubject}{splitlinethickness};
draw splitline withpen pencircle scaled thi withcolor linecolor;
draw (0, OverlayHeight - .5thi) withcolor background;
\stopuseMPgraphic
\stopenvironment
